% !TEX output_directory = ../publica
\documentclass[twoside, 10pt, a4paper]{article}
\usepackage[utf8]{inputenc}
\usepackage[T1]{fontenc}
\usepackage[portuguese]{babel}

% --- FONTE PADRÃO (Estilo Didático Encorpado) ---
\usepackage{helvet} 
\renewcommand{\familydefault}{\sfdefault}

% --- GEOMETRIA E ESPAÇAMENTO ---
\usepackage[top=1.6cm, bottom=1.5cm, left=0.9cm, right=0.9cm, headheight=22pt, headsep=0.4cm]{geometry}
\usepackage{microtype} 
\usepackage{setspace}
\setstretch{1.0}
\usepackage{parskip}
\setlength{\parskip}{2pt}
\setlength{\parindent}{0pt}

\newlength{\espacoPosTitulo}\setlength{\espacoPosTitulo}{0.3cm}
\newlength{\espacoPreQuestao}\setlength{\espacoPreQuestao}{0.1cm}
\newlength{\espacoQuestao}\setlength{\espacoQuestao}{0.2cm}
\newlength{\espacoAlt}\setlength{\espacoAlt}{0.2cm}

\usepackage{enumitem}
\setlist{itemsep=0.4cm, topsep=0.1cm, parsep=0pt, partopsep=0pt}
\newcommand{\larguraTikz}{0.85\linewidth} 

% --- MATEMÁTICA E CORES ---
\usepackage{amsmath, amsfonts, amssymb, nccmath, mathastext}
\usepackage{xcolor}
\definecolor{indigo}{HTML}{4F46E5}
\definecolor{CorTitUm}{HTML}{FF6F61}
\definecolor{CorTitDois}{HTML}{FF6F61}
\definecolor{CorTitTres}{HTML}{658894}
\definecolor{CorLinha}{HTML}{B0B0B0} 
\definecolor{metal}{HTML}{7F8C8D}
\definecolor{vidro}{HTML}{3498DB}
\definecolor{ouro}{HTML}{F1C40F}
\definecolor{concreto}{HTML}{95A5A6}
\definecolor{madeira}{HTML}{D35400}
\definecolor{CinzaEscuro}{HTML}{4A4A4A} 
\definecolor{CinzaMedio}{HTML}{848484} 
\definecolor{Cinzablack}{HTML}{353535} 

% --- MINI-CABEÇALHO E RODAPÉ AUTORAL (TWOSIDE) ---
\usepackage{fancyhdr}
\pagestyle{fancy}
\fancyhf{} 
\renewcommand{\headrulewidth}{0.3pt} 
\renewcommand{\footrulewidth}{0.3pt}
\fancyhead[LE]{\small\sffamily\bfseries\color{gray!75} ENERGIA MECÂNICA E POTÊNCIA}
\fancyhead[RO]{\small\sffamily\color{gray!75} \texttt{www.profraphaelpaulino.com.br}}
\fancyfoot[LE]{\small \textbf{\thepage}}
\fancyfoot[RO]{\small \textbf{\thepage}}

% --- CAIXAS DE DESTAQUE ---
\usepackage[most]{tcolorbox}
\newtcolorbox{boxresponda}{colback=white, colframe=gray!45, boxrule=0pt, leftrule=1.7pt, sharp corners, enlarge left by=0.4cm, width=\linewidth-0.4cm, left=8pt, right=0pt, top=2pt, bottom=2pt, fontupper=\small}
\definecolor{CorDicaBorda}{HTML}{17c1be} 
\definecolor{CorDicaFundo}{HTML}{f3fbfb} 
\newtcolorbox{boxdica}{colback=CorDicaFundo, colframe=CorDicaBorda, boxrule=0pt, leftrule=2.5pt, sharp corners, width=\linewidth, left=8pt, right=8pt, top=6pt, bottom=6pt, halign=center, fontupper=\small}

% --- TÍTULOS ---
\usepackage{titlesec}
\titleformat{\section}{\color{black}\large\bfseries}{\thesection}{0em}{\vspace{0.1cm}\titlerule[0.8pt]}
\titlespacing*{\section}{0pt}{10pt}{6pt} 
\titleformat{\subsubsection}[runin]{\color{black}\bfseries}{}{0em}{}
\titlespacing*{\subsubsection}{0pt}{0pt}{0.15cm}

% --- COLUNAS E GRÁFICOS ---
\usepackage{multicol}
\setlength{\columnsep}{1cm}
\setlength{\columnseprule}{0.5pt}
\def\columnseprulecolor{\color{CorLinha}}
\usepackage{adjustbox, tikz, pgfplots, circuitikz}
\usetikzlibrary{shadows, shapes.misc, positioning, calc}
\pgfplotsset{compat=1.18}

\begin{document}
\thispagestyle{empty}
\color{Cinzablack}
\vspace*{-1.8cm}

% --- CABEÇALHO PRINCIPAL AUTORAL (Apenas 1ª página) ---
\noindent
\begin{minipage}{\textwidth}
    \fontfamily{phv}\selectfont
    \noindent
    \begin{minipage}[t]{0.70\linewidth}
        {\Large \bfseries \MakeUppercase{Caderno de Atividades Suplementares}}\\[0.1cm]
        {\large \textmd{\textcolor{CinzaEscuro}{Unidade: ENERGIA MECÂNICA E POTÊNCIA}}}
    \end{minipage}%
    \begin{minipage}[t]{0.30\linewidth}
        \raggedleft
        \small \textbf{Prof. Raphael Paulino}\\
        \textsf{\small \textcolor{indigo}{\texttt{profraphaelpaulino.com.br}}}
    \end{minipage}
    
    \vspace{0.15cm}
    \noindent\rule{\linewidth}{1.2pt}
    \vspace{-0.1cm}
    \footnotesize\textsf{\textbf{ÁREA:} FÍSICA \hfill \textbf{NÍVEL:} EF II (9º Ano)}
    \vspace{0.1cm}
    \noindent\rule{\linewidth}{0.4pt}
\end{minipage}

\par\vspace{0.3cm} 
\raggedcolumns
\begin{multicols*}{2}
\vspace{0.1cm}

\subsection*{Capítulo 1: Energia Cinética}
\vspace{-0.2cm}\noindent\rule{\linewidth}{1.5pt} 
\par\vspace{\espacoPosTitulo}
\subsubsection*{1.} A figura ilustra um teste de desempenho automotivo em uma pista retilínea. Com base exclusivamente nos dados operacionais exibidos na telemetria, equacione literalmente o Teorema da Energia Cinética e determine o trabalho resultante realizado pelo motor do veículo entre os pontos \textbf{A} e \textbf{B}.
\begin{center}
% ==========================================
% CONTROLE INDIVIDUAL DESTA IMAGEM:
% 1. CORTAR BORDAS: Mude os zeros do trim (Esquerda, Baixo, Direita, Topo). Ex: trim=1cm 0cm 1cm 0cm
% 2. TAMANHO: Para encolher só esta imagem, troque \larguraTikz por um valor (Ex: 0.5\linewidth)
% ==========================================
\begin{adjustbox}{trim=0cm 0cm 0cm 0cm, clip, max width=\larguraTikz}
\begin{tikzpicture}
% --- CORES CUSTOMIZADAS ---
\definecolor{asfalto}{HTML}{2F3640}
\definecolor{ceu}{HTML}{D9EEF7}
\definecolor{carroceria}{HTML}{E74C3C}
\definecolor{roda}{HTML}{192A56}

% LAYER 1: Fundo e Pista
\fill[ceu, rounded corners=2pt] (-0.5,-1.5) rectangle (8.5, 3.5);
\fill[asfalto, drop shadow={opacity=0.15, shadow xshift=2pt, shadow yshift=-2pt}] (-0.5,-1.5) rectangle (8.5, 0);
\draw[dashed, white, thick] (-0.5, -0.75) -- (8.5, -0.75);

% LAYER 2: Veículo A
\begin{scope}[shift={(0.5, 0)}]
    \fill[carroceria, rounded corners=3pt, drop shadow={opacity=0.2}] (0,0.2) rectangle (2, 0.8);
    \fill[carroceria, rounded corners=2pt] (0.3,0.8) -- (0.6,1.3) -- (1.5,1.3) -- (1.8,0.8) -- cycle;
    \fill[vidro] (0.4,0.8) -- (0.65,1.2) -- (1.1,1.2) -- (1.1,0.8) -- cycle;
    \fill[vidro] (1.15,0.8) -- (1.15,1.2) -- (1.45,1.2) -- (1.7,0.8) -- cycle;
    \fill[roda] (0.4,0.2) circle (0.25);
    \fill[roda] (1.6,0.2) circle (0.25);
    \draw[->, ultra thick, white] (2.2, 0.6) -- (3.2, 0.6) node[above, font=\bfseries] {v\textsubscript{A} = 10 m/s};
    \node[white, font=\bfseries\small] at (1, 0.5) {m = 1200 kg};
\end{scope}

% LAYER 3: Veículo B
\begin{scope}[shift={(5.5, 0)}]
    \fill[carroceria, rounded corners=3pt, drop shadow={opacity=0.2}, opacity=0.5] (0,0.2) rectangle (2, 0.8);
    \fill[carroceria, rounded corners=2pt, opacity=0.5] (0.3,0.8) -- (0.6,1.3) -- (1.5,1.3) -- (1.8,0.8) -- cycle;
    \fill[roda, opacity=0.5] (0.4,0.2) circle (0.25);
    \fill[roda, opacity=0.5] (1.6,0.2) circle (0.25);
    \draw[->, ultra thick, white] (2.2, 0.6) -- (4.0, 0.6) node[above, font=\bfseries] {v\textsubscript{B} = 20 m/s};
\end{scope}

% LAYER 4: Marcadores
\draw[thick, white] (1.5, -0.2) -- (1.5, 0) node[below=0.2cm, fill=asfalto, inner sep=2pt] {Posição A};
\draw[thick, white] (6.5, -0.2) -- (6.5, 0) node[below=0.2cm, fill=asfalto, inner sep=2pt] {Posição B};
\end{tikzpicture}
\end{adjustbox}
\end{center}
\par\vspace{\espacoQuestao}

\noindent\textcolor{CorLinha}{\rule{\linewidth}{0.5pt}} 
\par\vspace{\espacoPreQuestao}
\subsubsection*{2.} A energia cinética é a grandeza escalar associada ao estado de movimento de um corpo. Julgue os itens a seguir como Verdadeiros (V) ou Falsos (F), justificando algebricamente as afirmativas consideradas falsas.
\begin{enumerate}[label=\Roman*., itemsep=\espacoAlt]
    \item Se a velocidade de um objeto for duplicada, sua energia cinética também será duplicada, pois são grandezas diretamente proporcionais.
    \item O trabalho total realizado pelas forças que atuam sobre uma partícula é igual à variação da sua energia cinética.
    \item É possível que um sistema apresente energia cinética negativa se o corpo estiver se movendo no sentido oposto ao referencial adotado.
\end{enumerate}
\par\vspace{\espacoQuestao}

\noindent\textcolor{CorLinha}{\rule{\linewidth}{0.5pt}} 
\par\vspace{\espacoPreQuestao}
\subsubsection*{3.} O gráfico ilustra o comportamento da força resultante \textbf{F} que atua sobre um bloco de massa \textbf{4 kg}, inicialmente em repouso na origem das posições, em função do deslocamento \textbf{d}. Extraia as informações geométricas do diagrama para determinar a velocidade do bloco ao atingir a marca de \textbf{10 m}.
\begin{center}
% ==========================================
% CONTROLE INDIVIDUAL DESTA IMAGEM:
% ==========================================
\begin{adjustbox}{trim=0cm 0cm 0cm 0cm, clip, max width=\larguraTikz}
\begin{tikzpicture}
\definecolor{grafico}{HTML}{27AE60}
\definecolor{fundo}{HTML}{F9F9F9}
\begin{axis}[
    width=7cm, height=5cm,
    axis lines=middle,
    xmin=0, xmax=12,
    ymin=0, ymax=60,
    xlabel={\textbf{d (m)}},
    ylabel={\textbf{F (N)}},
    grid=both,
    grid style={line width=.1pt, draw=gray!20},
    major grid style={line width=.2pt,draw=gray!50},
    ticklabel style={font=\footnotesize, fill=fundo, inner sep=1pt},
    enlargelimits={abs=0.5},
    axis background/.style={fill=fundo}
]
% --- APLICANDO DROP SHADOW NO GRÁFICO (ÁREA) ---
\addplot[fill=grafico, fill opacity=0.3, draw=grafico, thick] coordinates {(0,50) (8,50) (10,0)} \closedcycle;
\addplot[mark=*, mark size=1.5pt, grafico] coordinates {(0,50) (8,50) (10,0)};
\node[fill=fundo, inner sep=2pt] at (axis cs:4,25) {\textbf{Área = Trabalho}};
\end{axis}
\end{tikzpicture}
\end{adjustbox}
\end{center}
\par\vspace{\espacoQuestao}

\noindent\textcolor{CorLinha}{\rule{\linewidth}{0.5pt}}
\par\vspace{\espacoPreQuestao}
\subsubsection*{4.} Um estudante tentou resolver o seguinte problema: "Um projétil de massa \textbf{2 kg} impacta um anteparo a \textbf{10 m/s} e para completamente após perfurar \textbf{2 m}. Qual o módulo da força dissipativa constante?". A resolução do aluno foi:

\textit{Passo 1: Calculando a Energia: $E_c = m \cdot v = 2 \cdot 10 = 20 \text{ J}$.}

\textit{Passo 2: Usando o Teorema: $W = F \cdot d \rightarrow 20 = F \cdot 2 \rightarrow F = 10 \text{ N}$.}
\begin{boxresponda}
\textbf{Responda:} Diagnostique o erro conceitual cometido pelo estudante fictício no Passo 1, explique a regra matemática correta e recalcule o valor real da força resistiva.
\end{boxresponda}
\par\vspace{\espacoQuestao}

\noindent\textcolor{CorLinha}{\rule{\linewidth}{0.5pt}} 
\par\vspace{\espacoPreQuestao}
\subsubsection*{5.} (ENEM - Adaptada) O Conselho Nacional de Trânsito adverte constantemente sobre os riscos do excesso de velocidade. A severidade dos danos estruturais em uma colisão frontal é diretamente proporcional à energia cinética dissipada no instante do impacto. Considere dois veículos idênticos trafegando em uma mesma via. O veículo A desloca-se a \textbf{40 km/h}, enquanto o veículo B encontra-se a \textbf{120 km/h}. Com base no princípio físico da energia cinética, se ambos colidirem contra uma barreira idêntica, a energia dissipada pelo veículo B será:
\begin{enumerate}[label=\alph*), itemsep=\espacoAlt]
    \item equivalente ao triplo da energia dissipada pelo veículo A.
    \item equivalente ao sêxtuplo da energia dissipada pelo veículo A.
    \item equivalente a nove vezes a energia dissipada pelo veículo A.
    \item equivalente a doze vezes a energia dissipada pelo veículo A.
    \item igual à do veículo A, pois as massas dos veículos são idênticas.
\end{enumerate}
\begin{boxresponda}
\textbf{Responda:} Refute a alternativa "a", provando matematicamente por que a proporção direta sugerida por ela é um distrator incorreto.
\end{boxresponda}
\par\vspace{\espacoQuestao}

\noindent\textcolor{CorLinha}{\rule{\linewidth}{0.5pt}} 
\par\vspace{\espacoPreQuestao}
\subsubsection*{6.} (Estilo UNESP) O sistema de frenagem de emergência de uma composição ferroviária de \textbf{200 toneladas} é acionado quando o trem atinge a velocidade de \textbf{72 km/h}. Os freios aplicam uma força de atrito constante que realiza um trabalho total dissipativo para imobilizar a composição de forma segura. Aplique o Teorema da Energia Cinética, estruturando a equação algébrica passo a passo, para descobrir o módulo do trabalho mecânico realizado durante a frenagem.
\par\vspace{\espacoQuestao}

\subsection*{Capítulo 2: Energia Potencial}
\vspace{-0.2cm}\noindent\rule{\linewidth}{1.5pt} 
\par\vspace{\espacoPosTitulo}
\subsubsection*{7.} O arranjo logístico apresenta caixas armazenadas em um galpão de distribuição. A análise estrutural exige o cálculo das energias acumuladas no sistema. Considere a aceleração gravitacional no local como \textbf{10 m/s²}. Observe a posição relativa e a massa de cada volume na figura para calcular a diferença algébrica entre a energia potencial gravitacional da Caixa Alfa e da Caixa Beta, tomando o solo como referencial nulo.
\begin{center}
% ==========================================
% CONTROLE INDIVIDUAL DESTA IMAGEM:
% ==========================================
\begin{adjustbox}{trim=0cm 0cm 0cm 0cm, clip, max width=\larguraTikz}
\begin{tikzpicture}
\definecolor{parede}{HTML}{ECF0F1}
\definecolor{prateleira}{HTML}{7F8C8D}
\definecolor{caixaA}{HTML}{E67E22}
\definecolor{caixaB}{HTML}{9B59B6}
\definecolor{chao}{HTML}{BDC3C7}

% LAYER 1: Fundo e Solo
\fill[parede] (-1, 0) rectangle (7, 5);
\fill[chao, drop shadow={opacity=0.1}] (-1, -1) rectangle (7, 0);
\draw[thick] (-1,0) -- (7,0);

% LAYER 2: Prateleiras (com cantos arredondados)
\fill[prateleira, rounded corners=1pt] (0.5, 3.8) rectangle (3.5, 4.0);
\fill[prateleira, rounded corners=1pt] (3.5, 1.8) rectangle (6.5, 2.0);

% LAYER 3: Caixas (com sombras)
\fill[caixaA, rounded corners=2pt, drop shadow={opacity=0.2}] (1.0, 4.0) rectangle (2.5, 5.0);
\node[white, font=\bfseries] at (1.75, 4.5) {ALFA};
\node[white, font=\bfseries\scriptsize] at (1.75, 4.2) {20 kg};

\fill[caixaB, rounded corners=2pt, drop shadow={opacity=0.2}] (4.0, 2.0) rectangle (5.5, 3.0);
\node[white, font=\bfseries] at (4.75, 2.5) {BETA};
\node[white, font=\bfseries\scriptsize] at (4.75, 2.2) {50 kg};

% LAYER 4: Cotas de Altura
\draw[<->, thick, dashed] (0.2, 0) -- (0.2, 3.8) node[midway, fill=parede, inner sep=2pt, font=\bfseries] {h\textsubscript{1} = 4,0 m};
\draw[<->, thick, dashed] (6.8, 0) -- (6.8, 1.8) node[midway, fill=parede, inner sep=2pt, font=\bfseries] {h\textsubscript{2} = 1,5 m};
\end{tikzpicture}
\end{adjustbox}
\end{center}
\par\vspace{\espacoQuestao}

\noindent\textcolor{CorLinha}{\rule{\linewidth}{0.5pt}} 
\par\vspace{\espacoPreQuestao}
\subsubsection*{8.} Uma prensa mecânica industrial opera utilizando um sistema de amortecimento por molas. Um componente de massa \textbf{M} comprime a mola em uma distância \textbf{x}. Sabendo que a mola armazena uma energia potencial elástica \textbf{E\textsubscript{e}}, formule algebricamente a expressão da constante elástica \textbf{k} da mola em função da energia armazenada e da compressão sofrida, explicitando todos os passos lógicos.
\par\vspace{\espacoQuestao}

\noindent\textcolor{CorLinha}{\rule{\linewidth}{0.5pt}} 
\par\vspace{\espacoPreQuestao}
\subsubsection*{9.} Os engenheiros de teste plotaram o diagrama de deformação elástica para o feixe de molas de uma suspensão automotiva, que obedece perfeitamente à Lei de Hooke. Utilize os dados geométricos extraídos do gráfico para determinar a energia potencial elástica armazenada no sistema quando a deformação atingir \textbf{0,5 m}.
\begin{center}
% ==========================================
% CONTROLE INDIVIDUAL DESTA IMAGEM:
% ==========================================
\begin{adjustbox}{trim=0cm 0cm 0cm 0cm, clip, max width=\larguraTikz}
\begin{tikzpicture}
\definecolor{linhaelastica}{HTML}{8E44AD}
\definecolor{fundografico}{HTML}{FFFFFF}
\begin{axis}[
    width=7cm, height=5.5cm,
    axis lines=left,
    xmin=0, xmax=0.6,
    ymin=0, ymax=1200,
    xlabel={\textbf{Deformação x (m)}},
    ylabel={\textbf{Força Restauradora F (N)}},
    grid=major,
    grid style={dashed, gray!40},
    ticklabel style={font=\footnotesize},
    scaled y ticks=false,
]
% LAYER 1: Gráfico e Rótulos
\addplot[domain=0:0.4, linhaelastica, ultra thick] {2500*x};
\addplot[mark=*, mark size=2pt, linhaelastica] coordinates {(0.4, 1000)};
\draw[dashed, thick] (axis cs:0.4,0) -- (axis cs:0.4,1000) -- (axis cs:0,1000);
\node[anchor=south east] at (axis cs:0.4, 1000) {\textbf{Ponto de Teste}};
\end{axis}
\end{tikzpicture}
\end{adjustbox}
\end{center}
\par\vspace{\espacoQuestao}

\noindent\textcolor{CorLinha}{\rule{\linewidth}{0.5pt}} 
\par\vspace{\espacoPreQuestao}
\subsubsection*{10.} Sobre a natureza das energias potenciais, analise os itens e julgue-os como Verdadeiros (V) ou Falsos (F).
\begin{enumerate}[label=\Roman*., itemsep=\espacoAlt]
    \item A energia potencial gravitacional de um sistema independe do referencial adotado, sendo uma propriedade intrínseca da massa do corpo.
    \item Duas molas diferentes, com constantes elásticas \textbf{k\textsubscript{1}} e \textbf{k\textsubscript{2}}, sendo $k_1 > k_2$, se deformadas pela mesma distância, armazenarão quantidades diferentes de energia elástica.
    \item A energia potencial está sempre associada a forças conservativas, como o peso e a força elástica.
\end{enumerate}
\par\vspace{\espacoQuestao}
\noindent\textcolor{CorLinha}{\rule{\linewidth}{0.5pt}} 
\par\vspace{\espacoPreQuestao}
\subsubsection*{11.} Em um salto de \textit{Bungee Jump}, o praticante de massa \textbf{m} atinge o ponto mais baixo de sua trajetória, no qual a corda elástica, de constante \textbf{k}, apresenta um alongamento máximo \textbf{$\Delta x$}. Neste ponto de inversão do movimento, escreva a igualdade literal que relaciona a perda de energia potencial gravitacional do praticante com a energia potencial elástica armazenada na corda, desprezando a resistência do ar.
\par\vspace{\espacoQuestao}

\noindent\textcolor{CorLinha}{\rule{\linewidth}{0.5pt}} 
\par\vspace{\espacoPreQuestao}
\subsubsection*{12.} (FUVEST - Adaptada) O dimensionamento de turbinas em usinas hidrelétricas depende da carga hidráulica do reservatório. A água do reservatório flui pelos dutos e converte sua energia de posição em movimento. Considerando a densidade da água como \textbf{1.000 kg/m³} e a gravidade \textbf{g = 10 m/s²}, se um volume equivalente a \textbf{5 m³} despenca de um desnível de \textbf{120 m} de altura, determine a energia potencial gravitacional disponível nesse volume antes da queda.
\begin{enumerate}[label=\alph*), itemsep=\espacoAlt]
    \item \textbf{6.000 J}
    \item \textbf{60.000 J}
    \item \textbf{600.000 J}
    \item \textbf{6.000.000 J}
    \item \textbf{60.000.000 J}
\end{enumerate}
\begin{boxresponda}
\textbf{Responda:} Refute a alternativa "a", apresentando o cálculo correto da massa de água em quilogramas e provando matematicamente o erro de ordem de grandeza cometido nessa opção.
\end{boxresponda}
\par\vspace{\espacoQuestao}

\noindent\textcolor{CorLinha}{\rule{\linewidth}{0.5pt}} 
\par\vspace{\espacoPreQuestao}
\subsubsection*{13.} (Estilo UNESP) O diagrama ilustra um arqueiro de alta precisão esticando a corda de seu arco composto. A balança de tensão acoplada indica o comportamento linear das hastes flexíveis. Assumindo que a flecha possui massa de \textbf{40 g} ($ \textbf{0,04 kg} $) e que toda a energia mecânica do arco é transferida para o projétil sem perdas dissipativas, assinale a alternativa que indica a velocidade de lançamento da flecha assim que abandona a corda.
\begin{center}
% ==========================================
% CONTROLE INDIVIDUAL DESTA IMAGEM:
% 1. CORTAR BORDAS: Mude os zeros do trim (Esquerda, Baixo, Direita, Topo). Ex: trim=1cm 0cm 1cm 0cm
% 2. TAMANHO: Para encolher só esta imagem, troque \larguraTikz por um valor (Ex: 0.5\linewidth)
% ==========================================
\begin{adjustbox}{trim=0cm 0cm 0cm 0cm, clip, max width=\larguraTikz}
\begin{tikzpicture}
\definecolor{arco}{HTML}{D35400}
\definecolor{corda}{HTML}{2C3E50}
\definecolor{flecha}{HTML}{7F8C8D}
\definecolor{fundoarco}{HTML}{FDFEFE}

\fill[fundoarco, rounded corners=4pt, drop shadow={opacity=0.1}] (-2, -3) rectangle (6, 3);

% LAYER 1: Arco Tensionado
\draw[arco, line width=4pt] (0,2) to[bend left=30] (0,-2);
\draw[corda, thick] (0,2) -- (2.5,0) -- (0,-2);

% LAYER 2: Flecha
\draw[flecha, line width=2pt, ->] (2.5,0) -- (-1.5,0);
\fill[black] (-1.5,-0.1) -- (-1.8,0) -- (-1.5,0.1) -- cycle;

% LAYER 3: Cotas
\draw[<->, dashed, thick] (0, -0.5) -- (2.5, -0.5) node[midway, below, font=\bfseries] {Deformação = 0,8 m};
\node[fill=fundoarco, draw=black, inner sep=3pt, rounded corners=2pt, drop shadow={opacity=0.2}] at (3.5, 1.5) {\textbf{k = 500 N/m}};
\end{tikzpicture}
\end{adjustbox}
\end{center}
\begin{enumerate}[label=\alph*), itemsep=\espacoAlt]
    \item \textbf{20 m/s}
    \item \textbf{40 m/s}
    \item \textbf{60 m/s}
    \item \textbf{80 m/s}
    \item \textbf{100 m/s}
\end{enumerate}
\par\vspace{\espacoQuestao}

\noindent\textcolor{CorLinha}{\rule{\linewidth}{0.5pt}}
\par\vspace{\espacoPreQuestao}
\subsubsection*{14.} (ENEM - Adaptada) Pinballs antigos utilizam um lançador mecânico acionado por mola para colocar a esfera em jogo. Para modernizar o equipamento, propôs-se substituir a mola original por uma nova que possui o dobro da constante elástica ($ k_{nova} = 2k_{original} $), comprimindo exatamente a mesma distância \textbf{x}. Um estudante analisou a mudança e relatou o seguinte parecer:

\textit{Passo 1: A fórmula da energia potencial da mola é expressa por $ E_e = k \cdot x $.}

\textit{Passo 2: Como a constante dobrou e a distância \textbf{x} é a mesma, a energia elástica ficará multiplicada ao quadrado, ou seja, quatro vezes maior.}
\begin{boxresponda}
\textbf{Responda:} Diagnostique os dois erros cometidos pelo estudante fictício (um erro na fórmula descrita no Passo 1 e um erro de dedução lógica no Passo 2) e apresente qual será a verdadeira proporção de aumento da energia.
\end{boxresponda}
\par\vspace{\espacoQuestao}


\subsection*{Capítulo 3: Conservação de Energia Mecânica}
\vspace{-0.2cm}\noindent\rule{\linewidth}{1.5pt} 
\par\vspace{\espacoPosTitulo}
\subsubsection*{15.} A figura ilustra o arranjo experimental de um pêndulo simples utilizado em laboratório. A massa pendular é liberada do repouso na posição \textbf{A}. O nível de referência para a energia potencial gravitacional está posicionado no ponto mais baixo da trajetória (posição \textbf{B}). Deduza literalmente, passo a passo, a equação que determina a velocidade da massa no ponto \textbf{B}, em função da aceleração da gravidade \textbf{g}, do comprimento \textbf{L} e do ângulo \textbf{$\theta$} estabelecido na imagem.
\begin{center}
% ==========================================
% CONTROLE INDIVIDUAL DESTA IMAGEM:
% ==========================================
\begin{adjustbox}{trim=0cm 0cm 0cm 0cm, clip, max width=\larguraTikz}
\begin{tikzpicture}
% --- APLICANDO DROP SHADOW E CORES ---
\definecolor{teto}{HTML}{34495E}
\definecolor{esfera}{HTML}{E74C3C}
\definecolor{linhas}{HTML}{7F8C8D}

% LAYER 1: Teto e Fundo
\fill[teto, drop shadow={opacity=0.2, shadow xshift=1pt, shadow yshift=-2pt}] (-2, 0) rectangle (2, 0.4);
\draw[linhas, dashed, thick] (0,0) -- (0,-4) node[below, font=\bfseries] {Nível Zero (B)};

% LAYER 2: Pêndulo na Posição A
\draw[thick] (0,0) -- (2.5, -3);
\fill[esfera, drop shadow={opacity=0.3}] (2.5, -3) circle (0.3);
\node[white, font=\bfseries] at (2.5, -3) {m};

% LAYER 3: Pêndulo na Posição B (Fantasma)
\draw[thick, dashed, opacity=0.5] (0,0) -- (0, -3.7);
\fill[esfera, opacity=0.3] (0, -3.7) circle (0.3);

% LAYER 4: Cotas e Ângulos
\draw[->, thick, linhas] (0, -1) arc (270:309:1) node[midway, below=0.1cm, font=\bfseries] {$\theta$};
\node[fill=white, inner sep=2pt, font=\bfseries] at (1.25, -1.5) {L};
\draw[<->, dashed, thick] (3.0, -3.7) -- (3.0, -3.0) node[midway, right, font=\bfseries] {h};
\draw[dashed, thick] (0, -3.7) -- (3.2, -3.7);
\draw[dashed, thick] (2.5, -3.0) -- (3.2, -3.0);
\node[font=\bfseries, fill=white, inner sep=2pt] at (2.9, -2.5) {A};
\end{tikzpicture}
\end{adjustbox}
\end{center}
\par\vspace{\espacoQuestao}

\noindent\textcolor{CorLinha}{\rule{\linewidth}{0.5pt}} 
\par\vspace{\espacoPreQuestao}
\subsubsection*{16.} Em um sistema conservativo ideal, onde atuam exclusivamente as forças peso e elástica, julgue os itens a seguir como Verdadeiros (V) ou Falsos (F). Justifique obrigatoriamente as alternativas falsas utilizando o princípio da conservação.
\begin{enumerate}[label=\Roman*., itemsep=\espacoAlt]
    \item Se um corpo em queda livre diminui sua energia potencial gravitacional em \textbf{50 J}, sua energia cinética deve, simultaneamente, aumentar em \textbf{50 J}.
    \item A energia mecânica total de um sistema conservativo é nula em todos os pontos da trajetória, pois a soma algébrica das variações de energia é zero.
    \item A presença de atrito do ar transformaria parte da energia mecânica em energia térmica, caracterizando o sistema como não conservativo.
\end{enumerate}
\par\vspace{\espacoQuestao}

\noindent\textcolor{CorLinha}{\rule{\linewidth}{0.5pt}} 
\par\vspace{\espacoPreQuestao}
\subsubsection*{17.} Um bloco de \textbf{5 kg} desliza por uma pista com perfil curvo, cujas superfícies são perfeitamente polidas (sem atrito). O objeto inicia seu movimento a partir do repouso no ponto \textbf{P}. Extraia os desníveis topográficos apresentados na figura e calcule a velocidade do bloco ao atingir o platô no ponto \textbf{Q}. Considere $ \textbf{g} = \textbf{10 m/s²} $.
\begin{center}
% ==========================================
% CONTROLE INDIVIDUAL DESTA IMAGEM:
% ==========================================
\begin{adjustbox}{trim=0cm 0cm 0cm 0cm, clip, max width=\larguraTikz}
\begin{tikzpicture}
\definecolor{pista}{HTML}{95A5A6}
\definecolor{bloco}{HTML}{F39C12}
\definecolor{solo}{HTML}{ECF0F1}

% LAYER 1: Pista e Solo
\fill[solo] (-1, -1) rectangle (8, 0);
\draw[thick] (-1,0) -- (8,0);

% Construção da Pista com sombra
\fill[pista, drop shadow={opacity=0.15}] 
    (0,4) to[out=-70, in=180] (3,1) to[out=0, in=180] (5,2) -- (7,2) -- (7,0) -- (0,0) -- cycle;

% LAYER 2: Bloco no Ponto P
\fill[bloco, rounded corners=1pt, drop shadow={opacity=0.2}] (-0.2, 4) rectangle (0.2, 4.4);
\node[above, font=\bfseries] at (0, 4.4) {P};

% LAYER 3: Bloco no Ponto Q (Fantasma)
\fill[bloco, rounded corners=1pt, opacity=0.4] (5.8, 2) rectangle (6.2, 2.4);
\node[above, font=\bfseries] at (6.0, 2.4) {Q};

% LAYER 4: Cotas de Altura (Teste Cego)
\draw[<->, dashed, thick] (-0.5, 0) -- (-0.5, 4) node[midway, fill=white, inner sep=2pt, font=\bfseries] {12,0 m};
\draw[<->, dashed, thick] (7.5, 0) -- (7.5, 2) node[midway, fill=white, inner sep=2pt, font=\bfseries] {3,2 m};
\end{tikzpicture}
\end{adjustbox}
\end{center}
\par\vspace{\espacoQuestao}

\noindent\textcolor{CorLinha}{\rule{\linewidth}{0.5pt}} 
\par\vspace{\espacoPreQuestao}
\subsubsection*{18.} 
Um objeto de massa \textbf{m} é abandonado do alto de um edifício de altura \textbf{H}. Durante a queda, devido à resistência do ar, uma quantidade constante de energia mecânica é dissipada sob a forma de calor. Ao tocar o solo, a velocidade de impacto do objeto é \textbf{v}. Desenvolva a expressão literal que determina o trabalho realizado pela força de resistência do ar em função de \textbf{m}, \textbf{g}, \textbf{H} e \textbf{v}.
\begin{boxdica}
\textbf{Lembre-se:} O trabalho mecânico de forças dissipativas, como o atrito e a resistência do ar, equivale sempre à variação total da Energia Mecânica do sistema ($E_{final} - E_{inicial}$).
\end{boxdica}

\par\vspace{\espacoQuestao}

\noindent\textcolor{CorLinha}{\rule{\linewidth}{0.5pt}} 
\par\vspace{\espacoPreQuestao}
\subsubsection*{19.} O departamento de engenharia monitorou a descida de uma composição em uma montanha-russa ideal (sem atrito). O gráfico gerado pelo software de telemetria relaciona a variação das energias cinética (\textbf{E\textsubscript{c}}) e potencial (\textbf{E\textsubscript{p}}) em função do tempo \textbf{t}. A partir da leitura visual das intersecções, identifique o valor da Energia Mecânica Total do sistema e comprove numericamente que ela se mantém constante nos instantes $ \textbf{t} = \textbf{0 s} $ e $ \textbf{t} = \textbf{4 s} $.
\begin{center}
% ==========================================
% CONTROLE INDIVIDUAL DESTA IMAGEM:
% ==========================================
\begin{adjustbox}{trim=0cm 0cm 0cm 0cm, clip, max width=\larguraTikz}
\begin{tikzpicture}
\definecolor{ecin}{HTML}{2ECC71}
\definecolor{epot}{HTML}{3498DB}
\definecolor{fundo}{HTML}{FDFEFE}
\begin{axis}[
    width=7cm, height=5.5cm,
    axis lines=left,
    xmin=0, xmax=6,
    ymin=0, ymax=600,
    xlabel={\textbf{Tempo t (s)}},
    ylabel={\textbf{Energia (J)}},
    grid=major,
    grid style={dashed, gray!30},
    ticklabel style={font=\footnotesize},
    legend style={at={(0.5,1.05)}, anchor=south, legend columns=2, draw=none, fill=fundo},
]
% LAYER 1: Curvas de Energia
\addplot[domain=0:5, ecin, ultra thick, drop shadow={opacity=0.1}] {500 - 500*exp(-x)};
\addlegendentry{\textbf{E\textsubscript{c}}}

\addplot[domain=0:5, epot, ultra thick, drop shadow={opacity=0.1}] {500*exp(-x)};
\addlegendentry{\textbf{E\textsubscript{p}}}

% LAYER 2: Pontos Notáveis
\addplot[mark=*, mark size=2pt, ecin] coordinates {(0, 0) (4, 490)};
\addplot[mark=*, mark size=2pt, epot] coordinates {(0, 500) (4, 10)};
\end{axis}
\end{tikzpicture}
\end{adjustbox}
\end{center}
\par\vspace{\espacoQuestao}

\noindent\textcolor{CorLinha}{\rule{\linewidth}{0.5pt}} 
\par\vspace{\espacoPreQuestao}
\subsubsection*{20.} Um projétil é disparado verticalmente para cima com uma velocidade inicial \textbf{v\textsubscript{0}}. Desprezando a resistência do ar e adotando o ponto de disparo como referencial nulo de altura, em qual fração exata da altura máxima alcançada pelo projétil a sua energia cinética corresponderá à metade de sua energia potencial gravitacional? Demonstre a resolução utilizando o Teorema da Conservação e o equacionamento literal de frações, através da notação $ \mfrac{E_p}{2} $.
\par\vspace{\espacoQuestao}
\noindent\textcolor{CorLinha}{\rule{\linewidth}{0.5pt}} 
\par\vspace{\espacoPreQuestao}
\subsubsection*{21.} Em competições de saltos ornamentais, o atleta caminha até a extremidade do trampolim flexível, deprime a prancha para baixo e é impulsionado para o ar antes de mergulhar na piscina. Elabore um diagrama descritivo detalhando as transformações sucessivas de energia (Elástica, Cinética e Gravitacional) desde o instante de máxima deformação da prancha até o momento exato em que o atleta atinge a superfície da água.
\par\vspace{\espacoQuestao}

\noindent\textcolor{CorLinha}{\rule{\linewidth}{0.5pt}} 
\par\vspace{\espacoPreQuestao}
\subsubsection*{22.} No pátio de triagem, um bloco de \textbf{2 kg} é pressionado contra uma mola, contraindo-a em \textbf{0,5 m}. Ao ser liberado, o bloco percorre um trecho horizontal liso e sobe uma rampa rugosa, parando instantaneamente no ponto \textbf{S}. Determine a quantidade de energia mecânica transformada em calor pelo atrito durante a subida da rampa. Adote \textbf{g = 10 m/s²}.
\begin{center}
% ==========================================
% CONTROLE INDIVIDUAL DESTA IMAGEM:
% 1. CORTAR BORDAS: Mude os zeros do trim (Esquerda, Baixo, Direita, Topo). Ex: trim=1cm 0cm 1cm 0cm
% 2. TAMANHO: Para encolher só esta imagem, troque \larguraTikz por um valor (Ex: 0.5\linewidth)
% ==========================================
\begin{adjustbox}{trim=0cm 0cm 0cm 0cm, clip, max width=\larguraTikz}
\begin{tikzpicture}
\definecolor{base}{HTML}{7F8C8D}
\definecolor{mola}{HTML}{2C3E50}
\definecolor{bloco}{HTML}{16A085}

% LAYER 1: Superfícies
\fill[base, opacity=0.3] (-2, -0.5) rectangle (3, 0);
\draw[thick] (-2,0) -- (3,0);
\fill[base, opacity=0.5] (3, 0) -- (7, 2) -- (7, -0.5) -- (3, -0.5) -- cycle;
\draw[thick] (3,0) -- (7,2);

% LAYER 2: Parede e Mola
\fill[base] (-2, 0) rectangle (-1.5, 1.5);
\draw[mola, thick, decorate, decoration={zigzag, segment length=4pt, amplitude=3pt}] (-1.5, 0.5) -- (0, 0.5);
\node[fill=white, inner sep=1pt, font=\bfseries\scriptsize] at (-0.75, 1.0) {k = 400 N/m};

% LAYER 3: Bloco Inicial
\fill[bloco, rounded corners=1pt, drop shadow={opacity=0.2}] (0, 0) rectangle (1, 1);

% LAYER 4: Bloco Final
\begin{scope}[shift={(5, 1)}, rotate=26.5]
    \fill[bloco, rounded corners=1pt, opacity=0.5] (0, 0) rectangle (1, 1);
    \node[above, font=\bfseries] at (0.5, 1.0) {S};
\end{scope}

% LAYER 5: Cotas
\draw[<->, dashed, thick] (5.5, 0) -- (5.5, 1.25) node[midway, right, font=\bfseries] {h = 1,2 m};
\node[font=\bfseries, fill=white, inner sep=2pt] at (5.0, 0.5) {$\mu \neq 0$};
\end{tikzpicture}
\end{adjustbox}
\end{center}
\par\vspace{\espacoQuestao}

\noindent\textcolor{CorLinha}{\rule{\linewidth}{0.5pt}} 
\par\vspace{\espacoPreQuestao}
\subsubsection*{23.} (ENEM - Adaptada) Em uma simulação de segurança viária, um automóvel de massa \textbf{m} desloca-se a uma velocidade de \textbf{108 km/h} (ou seja, \textbf{30 m/s}). O motorista aciona os freios abruptamente, travando as quatro rodas. O veículo derrapa e percorre uma distância retilínea até a imobilização completa. Considerando que a energia térmica gerada no asfalto equivale ao módulo do trabalho da força de atrito, deduza a expressão literal para o coeficiente de atrito cinético \textbf{$\mu$} em função de \textbf{v}, \textbf{g} e da distância de parada \textbf{d}.
\par\vspace{\espacoQuestao}

\noindent\textcolor{CorLinha}{\rule{\linewidth}{0.5pt}} 
\par\vspace{\espacoPreQuestao}
\subsubsection*{24.} (Estilo FUVEST) O diagrama mostra uma pista de acrobacias contendo um looping circular. Para que uma esfera rígida de massa \textbf{m} complete a volta inteira sem perder contato com a pista no ponto mais alto (\textbf{T}), ela deve possuir uma velocidade mínima crítica. Analise a geometria e estabeleça, rigorosamente, a equação algébrica que define a altura mínima de lançamento \textbf{h} em função exclusivamente do raio \textbf{R} da trajetória circular. Assuma que não há forças dissipativas e que a esfera parte do repouso.
\begin{center}
% ==========================================
% CONTROLE INDIVIDUAL DESTA IMAGEM:
% ==========================================
\begin{adjustbox}{trim=0cm 0cm 0cm 0cm, clip, max width=\larguraTikz}
\begin{tikzpicture}
\definecolor{metal}{HTML}{BDC3C7}
\definecolor{esfera}{HTML}{8E44AD}

% LAYER 1: Pista (Looping)
\draw[thick, double=metal, double distance=3pt] (-3, 4) to[out=-60, in=180] (0, 0) arc (-90:270:1.5) to[out=0, in=180] (5,0);

% LAYER 2: Esfera Inicial
\fill[esfera, drop shadow={opacity=0.3}] (-3, 4.3) circle (0.2);
\draw[dashed, thick] (-3, 0) -- (-3, 4.3);
\draw[<->, thick] (-3.5, 0) -- (-3.5, 4.3) node[midway, left, font=\bfseries] {h};

% LAYER 3: Looping Cotas
\draw[dashed, thick] (1.5, 1.5) -- (1.5, 3.0) node[midway, right, font=\bfseries] {R};
\fill[black] (1.5, 1.5) circle (1.5pt);
\node[above, font=\bfseries] at (1.5, 3.2) {T};
\end{tikzpicture}
\end{adjustbox}
\end{center}
\par\vspace{\espacoQuestao}

\noindent\textcolor{CorLinha}{\rule{\linewidth}{0.5pt}} 
\par\vspace{\espacoPreQuestao}
\subsubsection*{25.} (Estilo VUNESP) Um praticante de \textit{Snowboard} de \textbf{70 kg} inicia sua descida, a partir do repouso, do topo de uma encosta gelada. Ao longo da descida, devido à leve camada de neve irregular, exatos \textbf{15\%} da energia mecânica inicial do atleta é dissipada por atrito e resistência aerodinâmica. Se o ponto de partida encontra-se a \textbf{40 m} acima da base plana, determine a velocidade com que o praticante atinge a parte inferior da pista. Adote \textbf{g = 10 m/s²}.
\par\vspace{\espacoQuestao}

\noindent\textcolor{CorLinha}{\rule{\linewidth}{0.5pt}} 
\par\vspace{\espacoPreQuestao}
\subsubsection*{26.} (FUVEST - Adaptada) Uma bola de basquete é solta, a partir do repouso, de uma altura de \textbf{2,0 m} em relação a um piso rígido de concreto. A cada colisão com o solo, ocorre uma deformação inelástica que dissipa \textbf{20\%} da energia mecânica da bola sob forma de calor e som. Com base no princípio da conservação fracionária, qual será a altura máxima atingida pela bola imediatamente após o seu segundo quique no chão?
\begin{enumerate}[label=\alph*), itemsep=\espacoAlt]
    \item \textbf{1,28 m}
    \item \textbf{1,44 m}
    \item \textbf{1,60 m}
    \item \textbf{1,80 m}
    \item \textbf{1,92 m}
\end{enumerate}
\begin{boxresponda}
\textbf{Responda:} Prove matematicamente por que a alternativa "c" está incorreta, detalhando a falha na lógica de subtrair diretamente as porcentagens em vez de aplicar os descontos sucessivos a cada impacto.
\end{boxresponda}
\par\vspace{\espacoQuestao}


\subsection*{Capítulo 4: Potência e Rendimento}
\vspace{-0.2cm}\noindent\rule{\linewidth}{1.5pt} 
\par\vspace{\espacoPosTitulo}
\subsubsection*{27.} O gráfico a seguir descreve o comportamento da potência mecânica desenvolvida pelo motor de um guindaste em função do tempo durante o içamento de uma viga de aço. Determine o trabalho mecânico total realizado pelo motor nos primeiros \textbf{20 s} de operação.
\begin{center}
% ==========================================
% CONTROLE INDIVIDUAL DESTA IMAGEM:
% ==========================================
\begin{adjustbox}{trim=0cm 0cm 0cm 0cm, clip, max width=\larguraTikz}
\begin{tikzpicture}
\definecolor{grafico}{HTML}{E67E22}
\definecolor{fundo}{HTML}{FDFEFE}
\begin{axis}[
    width=7cm, height=5cm,
    axis lines=middle,
    xmin=0, xmax=25,
    ymin=0, ymax=6000,
    xlabel={\textbf{Tempo t (s)}},
    ylabel={\textbf{Potência P (W)}},
    grid=both,
    grid style={line width=.1pt, draw=gray!20},
    major grid style={line width=.2pt,draw=gray!50},
    ticklabel style={font=\footnotesize, fill=fundo, inner sep=1pt},
    enlargelimits={abs=0.5},
    axis background/.style={fill=fundo}
]
% --- APLICANDO DROP SHADOW NO GRÁFICO (ÁREA) ---
\addplot[fill=grafico, fill opacity=0.3, draw=grafico, thick] coordinates {(0,0) (5,5000) (15,5000) (20,0)} \closedcycle;
\addplot[mark=*, mark size=1.5pt, grafico] coordinates {(0,0) (5,5000) (15,5000) (20,0)};
\node[fill=fundo, inner sep=2pt] at (axis cs:10, 2500) {\textbf{Área = Trabalho}};
\end{axis}
\end{tikzpicture}
\end{adjustbox}
\end{center}
\par\vspace{\espacoQuestao}

\noindent\textcolor{CorLinha}{\rule{\linewidth}{0.5pt}} 
\par\vspace{\espacoPreQuestao}
\subsubsection*{28.} Um elevador de carga industrial com massa total de \textbf{800 kg} é içado verticalmente com velocidade constante de \textbf{2 m/s} por um guincho elétrico. Analise o arranjo físico apresentado e equacione literalmente a potência útil desenvolvida pelo motor ao tracionar o cabo de aço, explicitando todos os raciocínios e vetores. Considere \textbf{g = 10 m/s²}.
\begin{center}
% ==========================================
% CONTROLE INDIVIDUAL DESTA IMAGEM:
% ==========================================
\begin{adjustbox}{trim=0cm 0cm 0cm 0cm, clip, max width=\larguraTikz}
\begin{tikzpicture}
\definecolor{motor}{HTML}{2C3E50}
\definecolor{cabo}{HTML}{7F8C8D}
\definecolor{elevador}{HTML}{D35400}
\definecolor{teto}{HTML}{BDC3C7}

% LAYER 1: Estrutura Fixa
\fill[teto, drop shadow={opacity=0.2}] (-2, 4) rectangle (2, 4.5);
\fill[motor, rounded corners=2pt, drop shadow={opacity=0.3}] (-0.8, 4.5) rectangle (0.8, 5.2);
\node[white, font=\bfseries] at (0, 4.85) {MOTOR};
\draw[thick, white] (-0.3, 4.5) -- (-0.3, 4.0);

% LAYER 2: Cabos
\draw[cabo, ultra thick] (0, 4) -- (0, 1.5);

% LAYER 3: Elevador
\fill[elevador, rounded corners=2pt, drop shadow={opacity=0.2}] (-1, 0) rectangle (1, 1.5);
\draw[white, thick] (-0.8, 0.2) rectangle (0.8, 1.3);
\node[white, font=\bfseries] at (0, 0.75) {800 kg};

% LAYER 4: Vetores
\draw[->, ultra thick] (0, 1.7) -- (0, 2.7) node[right=0.1cm, font=\bfseries] {v = 2 m/s};
\draw[->, ultra thick, dashed] (0, 0) -- (0, -1.0) node[right=0.1cm, font=\bfseries] {Peso (P)};
\draw[->, ultra thick, dashed] (0, 1.5) -- (0, 2.5) node[left=0.1cm, font=\bfseries] {Tração (T)};
\end{tikzpicture}
\end{adjustbox}
\end{center}
\par\vspace{\espacoQuestao}

\noindent\textcolor{CorLinha}{\rule{\linewidth}{0.5pt}} 
\par\vspace{\espacoPreQuestao}
\subsubsection*{29.} Sobre os conceitos termodinâmicos e mecânicos de Potência e Rendimento, julgue os itens a seguir como Verdadeiros (V) ou Falsos (F).
\begin{enumerate}[label=\Roman*., itemsep=\espacoAlt]
    \item A potência é uma grandeza física que quantifica a rapidez com que o trabalho é realizado ou a energia é transferida em um sistema.
    \item O rendimento de uma máquina real pode atingir exatos \textbf{100\%} se as forças de atrito e outras resistências dissipativas forem reduzidas a zero absoluto no laboratório.
    \item Se dois motores realizam o mesmo trabalho mecânico, aquele que realizar a tarefa no menor intervalo de tempo desenvolverá a maior potência útil.
\end{enumerate}
\par\vspace{\espacoQuestao}

\noindent\textcolor{CorLinha}{\rule{\linewidth}{0.5pt}} 
\par\vspace{\espacoPreQuestao}
\subsubsection*{30.} O projeto de um veículo subcompacto indica que o seu motor a combustão apresenta uma potência total de \textbf{120 kW}. Ao calcular a potência útil efetivamente entregue pelo eixo, sabendo que o rendimento global é de \textbf{40\%}, um técnico em treinamento estruturou a seguinte resposta:

\textit{Passo 1: Se o rendimento é de 40\%, significa que o motor perde 40\% da energia para o calor.}

\textit{Passo 2: A potência dissipada é, portanto, $ 0,40 \cdot 120 = 48 \text{ kW} $.}

\textit{Passo 3: A potência útil que sobra para girar o eixo é $ 120 - 48 = 72 \text{ kW} $.}
\begin{boxresponda}
\textbf{Responda:} Diagnostique o erro de premissa cometido no Passo 1 pelo técnico, defina corretamente o que significa um rendimento de \textbf{40\%} em uma máquina térmica e refaça o cálculo determinando os valores verdadeiros das potências Útil e Dissipada.
\end{boxresponda}
\par\vspace{\espacoQuestao}

\noindent\textcolor{CorLinha}{\rule{\linewidth}{0.5pt}} 
\par\vspace{\espacoPreQuestao}
\subsubsection*{31.} (ENEM - Adaptada) Um sistema de irrigação rural utiliza uma motobomba de captação para retirar água de um poço artesiano e despejá-la em um reservatório superior. O equipamento possui um rendimento eletromecânico de \textbf{60\%}. O regime de trabalho contínuo da bomba extrai e eleva \textbf{300 litros} de água por minuto a uma altura vertical de \textbf{20 m}. Considerando a densidade da água como \textbf{1 kg/L} e \textbf{g = 10 m/s²}, a potência elétrica total consumida da rede elétrica pelo motor da bomba aproxima-se de:
\begin{enumerate}[label=\alph*), itemsep=\espacoAlt]
    \item \textbf{600 W}
    \item \textbf{1.000 W}
    \item \textbf{1.666 W}
    \item \textbf{3.000 W}
    \item \textbf{6.000 W}
\end{enumerate}
\par\vspace{\espacoQuestao}

\noindent\textcolor{CorLinha}{\rule{\linewidth}{0.5pt}} 
\par\vspace{\espacoPreQuestao}
\subsubsection*{32.} (Estilo UNESP) Uma esteira transportadora inclinada opera em uma planta de mineração. O perfil topográfico mostra que a estrutura eleva o minério através de um desnível constante \textbf{h}. A esteira é alimentada a uma taxa de \textbf{40 kg/s} e transporta o material a uma velocidade rigorosamente constante. Sabendo que o motor que traciona o conjunto desenvolve uma potência útil de \textbf{6.000 W} exclusivamente para vencer o trabalho gravitacional, determine, por meio de modelagem algébrica, a altura vertical \textbf{h} alcançada pela esteira em \textbf{1 segundo} de operação. Assuma \textbf{g = 10 m/s²}.
\begin{center}
% ==========================================
% CONTROLE INDIVIDUAL DESTA IMAGEM:
% ==========================================
\begin{adjustbox}{trim=0cm 0cm 0cm 0cm, clip, max width=\larguraTikz}
\begin{tikzpicture}
\definecolor{esteira}{HTML}{34495E}
\definecolor{minerio}{HTML}{7F8C8D}
\definecolor{chao}{HTML}{BDC3C7}
\definecolor{apoio}{HTML}{95A5A6}

% LAYER 1: Solo
\fill[chao] (-1, -1) rectangle (8, 0);
\draw[thick] (-1,0) -- (8,0);

% LAYER 2: Apoios da Esteira
\fill[apoio] (1, 0) rectangle (1.2, 1);
\fill[apoio] (5, 0) rectangle (5.2, 3);
\fill[apoio] (6.8, 0) rectangle (7.0, 3.9);

% LAYER 3: Esteira Inclinada
\draw[esteira, line width=6pt] (0, 0.5) -- (7, 4);
\fill[black] (0, 0.5) circle (0.15);
\fill[black] (7, 4) circle (0.15);

% LAYER 4: Blocos de Minério (Drop Shadow)
\foreach \x/\y in {1/1, 2.5/1.75, 4/2.5, 5.5/3.25} {
    \fill[minerio, rounded corners=1pt, drop shadow={opacity=0.3}] (\x, \y) rectangle (\x+0.8, \y+0.4);
}

% LAYER 5: Cotas
\draw[<->, dashed, thick] (7.5, 0) -- (7.5, 4) node[midway, right, font=\bfseries] {Desnível h};
\draw[dashed, thick] (7, 4) -- (7.8, 4);
\draw[->, thick, rotate=26.5] (2, 1.5) -- (4, 1.5) node[midway, above, font=\bfseries] {v = const};
\end{tikzpicture}
\end{adjustbox}
\end{center}
\par\vspace{\espacoQuestao}

\end{multicols*}
\end{document}
\end{multicols*}
\end{document}